\section{A relative error bound for finite spheres}
\label{sec:transfer}
The limiting Gram matrices have kernels, so their positivity cannot be
transferred by assuming that their ranges vary continuously.  Instead, we
compare the annihilators themselves.  An exact diagonal change of variables
reduces every error to a fixed finite family of pair-generated annihilators.
This gives an error relative to $H_Q$, uniformly on the full Fock space.

Geometry-dependent diagonal changes of occupation-number coordinates
already relate pseudopotential annihilators and their zero modes~\cite[Section~II.E]{OrtizNussinovDukelskySeidel}.
For a gap estimate, the needed additional statement is quantitative:
on a fixed set of modes the change of variables approaches the identity,
and its error is controlled by pair energy. The proof below establishes
that statement before adding spectator modes and averaging over rotations.

For the local argument, let $\FF_{\mathrm{loc}}$ be the exterior algebra
of $\mathbb C^{25}$ with orthonormal basis $e_0,\ldots,e_{24}$.
For $Q\ge24$ we identify it with the Fock space of the first 25 orbital
modes in $U_Q$. By \eqref{eq:pairstar}, each limiting pair vector $v_p^*$
with $p\le23$ lies in $\bigwedge^2\mathbb C^{25}$.
Contraction by these vectors is defined on $\FF_{\mathrm{loc}}$
by the same exterior-algebra convention as before.

\begin{lemma}[Compression and annihilator bounds]\label{lem:compression}
Let $W$ be a linear map between finite-dimensional Hilbert spaces,
$G=W^\dagger W$, and let $C$ be self-adjoint on the domain of $W$. Then
\[
 GCG\ge0\quad\Longleftrightarrow\quad WCW^\dagger\ge0.
\]
On $\FF_{\mathrm{loc}}$, for any fixed finite collection
$A_b=c_kc_jB(v_p^*)$ with $0\le p\le23$ and $0\le j,k\le24$,
\begin{equation}\label{eq:annihilatorbound}
 \sum_b\norm{A_b\psi}^2\le C_0\ip\psi{h_*\psi},
 \qquad h_*:=\sum_{p=0}^{23}B(v_p^*)^\dagger B(v_p^*),
\end{equation}
where $C_0$ depends only on the collection.
\end{lemma}
\begin{proof}
The two positivity statements both assert that the quadratic form of $C$
is nonnegative on $\ran G=\ran W^\dagger$.  For the second assertion, use
$\norm{c_kc_jB(v_p^*)\psi}\le\norm{B(v_p^*)\psi}$ and sum; one may take
$C_0$ to be the largest multiplicity of a pair label $p$.
\end{proof}

\Needspace{11\baselineskip}
\begin{lemma}[Uniform transfer]\label{lem:transfer}
For $Q\ge24$ and $1\le D\le23$, put
\[
 C_D^Q=(d_D/d_a)\chi-E_D^Q+\eps I,
 \qquad h_Q=\sum_{p=0}^{23}B_p^\dagger B_p.
\]
Assuming \eqref{eq:4certificate}, there is a scalar $\rho_Q\ge0$,
independent of $D$ and tending to zero as $Q\to\infty$, such that
\begin{equation}\label{eq:relative}
 \sum_{r,s\in\mathcal R_D}(C_D^Q)_{rs}
 B(w_{Dr}^Q)^\dagger B(w_{Ds}^Q)
 \ge-\rho_Qh_Q
\end{equation}
on the entire Fock space $\FF_Q$.
\end{lemma}
\begin{proof}
\emph{An exact change of variables.}
Take $Q\ge24$.  All annihilators in \eqref{eq:relative}, including those
in $h_Q$, involve only the modes $0,\ldots,24$. First work on
$\FF_{\mathrm{loc}}$. Define the positive diagonal operator
\[
 \Delta_Q\ket A=\left(\prod_{x\in A}d_Q(x)\right)\ket A,
 \qquad d_Q(x)=\sqrt{(Q)_x/Q^x},
\]
on each basis wedge $\ket A$.  This local Fock space is fixed as $Q$
varies.  Thus $\Delta_Q$ and $\Delta_Q^{-1}$ converge in norm to $I$;
also $0<\Delta_Q\le I$.

The pair coefficient formula \eqref{eq:pair} gives
\[
 v_p(x,y)=r_p(Q)d_Q(x)d_Q(y)v_p^*(x,y),
 \qquad r_p(Q)=\sqrt{(2Q)^p/(2Q-2)_p}\ge1.
\]
Checking the factors attached to the removed orbital labels on a basis
wedge gives the exact identities
\begin{align}
 B_p&=r_p(Q)\Delta_Q^{-1}B(v_p^*)\Delta_Q,
       \label{eq:pairconjugation}\\
 B(v_p\wedge e_j\wedge e_k)
 &=\frac{r_p(Q)}{d_Q(j)d_Q(k)}\,
   \Delta_Q^{-1}B(v_p^*\wedge e_j\wedge e_k)\Delta_Q.
       \label{eq:fourconjugation}
\end{align}
These are identities on the local Fock space, with the inverse diagonal
acting on the output particle sector.  Repeated removed labels make both
sides zero.  In particular, solving \eqref{eq:pairconjugation} for
$B(v_p^*)\Delta_Q$ and using $r_p(Q)\ge1$ and $\norm{\Delta_Q}\le1$
gives
\begin{equation}\label{eq:pairenergytransfer}
 \Delta_Q h_*\Delta_Q\le h_Q.
\end{equation}

\emph{The error is controlled by pair energy.}
Fix $D$.  Index a finite family by
$b=(p,j,k)$ with $p+j+k=D$ and $j<k$, and put
\[
 A_b=B(v_p^*\wedge e_j\wedge e_k),\qquad
 \beta_{rb}^Q=S_r^Q(D,D;p,j,k)
                \frac{r_p(Q)}{d_Q(j)d_Q(k)}.
\]
Equations \eqref{eq:wQ} and \eqref{eq:fourconjugation} yield
\[
 B(w_{Dr}^Q)=\Delta_Q^{-1}
             \left(\sum_b\beta_{rb}^QA_b\right)\Delta_Q.
\]
The coefficients are real, and Lemma~\ref{lem:coupling} gives
$\beta_{rb}^Q\to\beta_{rb}^*:=S_r(D,D;p,j,k)$.  Hence the finite
real symmetric matrices
\[
 M_Q=(\beta^Q)^TC_D^Q\beta^Q
 \quad\hbox{converge to}\quad
 M_*=(\beta^*)^T(2\chi-E_D+\eps I)\beta^*.
\]
The quadratic expression in \eqref{eq:relative} is consequently
\[
 \Delta_Q\left(\sum_{b,c}(M_Q)_{bc}
                 A_b^\dagger\Delta_Q^{-2}A_c\right)\Delta_Q.
\]
Its limiting expression before the outer diagonals is
\begin{equation}\label{eq:limitpositive}
 \sum_{b,c}(M_*)_{bc}A_b^\dagger A_c
 =\sum_{r,s}(2\chi-E_D+\eps I)_{rs}
       B(w_{Dr}^*)^\dagger B(w_{Ds}^*)\ge0.
\end{equation}
Indeed, the certificate \eqref{eq:4certificate} and
Lemma~\ref{lem:compression}, applied to the map with columns $w_{Dr}^*$,
give positivity before lifting; the lift $\LL_4$ preserves positivity.

For clarity, the perturbation estimate does not require $M_*\ge0$.
Let $R_Q$ be the block matrix with entries
\[
 (R_Q)_{bc}=(M_Q)_{bc}\Delta_Q^{-2}-(M_*)_{bc}I.
\]
Regard $R_Q$ as an operator on the Hilbert direct sum of one copy of
$\FF_{\mathrm{loc}}$ for each index $b$.  The diagonal $\Delta_Q^{-2}$
acts on the output of $A_c$, hence on the sector with four fewer particles;
using its direct sum over local sectors makes this a single operator
identity. Its norm tends to zero, since both its index set and the local
Fock space are fixed. For every vector $\phi$,
\[
 \left|\sum_{b,c}\ip{A_b\phi}{(R_Q)_{bc}A_c\phi}\right|
 \le\norm{R_Q}\sum_b\norm{A_b\phi}^2
 \le C_0\norm{R_Q}\ip\phi{h_*\phi}.
\]
Combine this estimate with \eqref{eq:limitpositive}, take
$\phi=\Delta_Q\psi$, and apply \eqref{eq:pairenergytransfer}.  This
proves \eqref{eq:relative} on the local Fock space with a coefficient
tending to zero.  Taking a maximum over the fixed set $1\le D\le23$
gives a common $\rho_Q$.

\emph{Additional particles are spectators.}
Order the 25 local modes before all other modes in the Fock
factorization.  Under the isometry that sends a local wedge tensored with a spectator
wedge to their ordered product, each $c_j$ with $j\le24$ is its local
annihilator tensored with the identity. The fermionic parity factor belongs
to annihilators in the later spectator modes, which do not occur here.
Thus every operator in \eqref{eq:relative} acts on the local factor alone.
Tensoring the local inequality with the identity proves
it on $\FF_Q$, with the same $\rho_Q$, regardless of the number of
occupied modes outside this factor.
\end{proof}

The allowance $\eps I$ in the certificate must also be paid for in pair
energy.  Here it is useful to sum over the orthonormal highest vectors
\emph{before} estimating their wedges.  Bessel's inequality then avoids
an extra factor equal to the number of spin copies.

\begin{proposition}[Four-body estimate]\label{prop:A4}
For the coefficients in Appendix~\ref{app:certificate},
\begin{equation}\label{eq:A4bound}
 A_4\le\LL_4(W_4W_4^\dagger)+(1222\eps+o(1))H,
 \qquad Q\to\infty,
\end{equation}
uniformly on $\FF_Q$.
\end{proposition}
\begin{proof}
Haar averaging the matrix unit
$\ket{D,D;r}_Q\bra{D,D;s}_Q$ gives the corresponding copy matrix unit
tensored with the spin identity, divided by $d_D$.  Therefore, on the
retained spin blocks, the target minus $A_4$, with the $\eps I$ allowance
added, is
\[
 d_a\sum_{D=1}^{23}\int\sum_{r,s}(C_D^Q)_{rs}
 B(w_{Dr}^Q(g))^\dagger B(w_{Ds}^Q(g))\,dg.
\]
Lemma~\ref{lem:transfer}, conjugated by each rotation, bounds this below
by $-o(1)H$, because Schur averaging gives
\[
 d_a\int h_Q(g)\,dg=24H.
\]
The number of retained blocks is fixed, so this error remains uniform
in particle number.

To bound the allowance, fix $D$ and write
\[
 \ket{D,D;r}_Q
   =\sum_{\substack{p+j+k=D\\j<k}}S_{rb}^Q\,
                     v_p\otimes(e_j\wedge e_k),
 \qquad b=(p,j,k).
\]
The uncoupled vectors on the right are orthonormal.  So are the highest
vectors on the left: distinct $r$ belong to orthogonal spin summands of
the second pair factor.  Thus $S^Q(S^Q)^T=I$.  Apply the contraction
$S^Q\otimes I$ to the Hilbert-space-valued vector
$(c_kc_jB_p\psi)_b$.  It follows that
\[
 \sum_{r\in\mathcal R_D}\norm{B(w_{Dr}^Q)\psi}^2
 \le\sum_{\substack{p+j+k=D\\j<k}}\norm{c_kc_jB_p\psi}^2
 \le\sum_{\substack{p+j+k=D\\j<k}}\norm{B_p\psi}^2.
\]
This uses orthogonality before wedging; the vectors $w_{Dr}^Q$ themselves
need not be orthogonal.  The number of triples on the right is
\[
 n_D=\sum_{n=1}^D\left\lceil\frac n2\right\rceil
     =\left\lfloor\frac{(D+1)^2}{4}\right\rfloor.
\]
Conjugate by rotations and integrate.  Each pair norm contributes
$\ip\psi{H\psi}/d_a$, so the total allowance is bounded by
\[
 \eps\sum_{D=1}^{23}n_DH
 =\eps\left(\sum_{m=1}^{12}m^2+
                 \sum_{m=1}^{11}m(m+1)\right)H
 =1222\eps H.
\]
The discarded target spin blocks are positive.  Removing the allowance
therefore proves \eqref{eq:A4bound}.
\end{proof}
